\documentclass[12pt,a4paper]{article}

\usepackage[margin=2.5cm]{geometry}
\usepackage{amsmath,amssymb}
\usepackage{booktabs}
\usepackage{graphicx}
\usepackage{hyperref}
\usepackage{url}
\usepackage{microtype}
\usepackage{setspace}
\usepackage{titlesec}
\usepackage{abstract}
\usepackage{natbib}
\usepackage{caption}
\usepackage{subcaption}
\usepackage{listings}
\usepackage{xcolor}
\usepackage{float}
\usepackage{array}

\hypersetup{
  colorlinks=true,
  linkcolor=blue!60!black,
  citecolor=blue!60!black,
  urlcolor=blue!60!black
}

\lstset{
  basicstyle=\ttfamily\small,
  backgroundcolor=\color{gray!8},
  frame=single,
  rulecolor=\color{gray!40},
  breaklines=true,
  captionpos=b
}

\titleformat{\section}{\large\bfseries}{\thesection.}{0.5em}{}
\titleformat{\subsection}{\normalsize\bfseries}{\thesubsection.}{0.5em}{}

\onehalfspacing

% ---------------------------------------------------------------
\title{%
  \textbf{Lunar Ice Intelligence:}\\
  Physics-Informed Deep Learning for Water Ice Detection\\
  in Lunar Permanently Shadowed Regions
}

\author{
  Yago Almeida da Silva\\
  \small ORCID: \href{https://orcid.org/0009-0007-0094-0915}{0009-0007-0094-0915}\\
  \small \texttt{yagoalcontact@gmail.com}
}

\date{May 2026 \\ \small Preprint — not peer-reviewed}
% ---------------------------------------------------------------

\begin{document}

\maketitle

\begin{abstract}
We present \textit{Lunar Ice Intelligence}, an open-source platform for detecting
water ice in lunar permanently shadowed regions (PSRs) using physics-informed deep
learning. The system integrates real data from four NASA Lunar Reconnaissance Orbiter
(LRO) instruments — Diviner, Mini-RF, LROC WAC, and LAMP — with a multimodal neural
network (LunarCNN) that couples a convolutional image encoder with a physics-based
subsurface thermal encoder. Labels are derived exclusively from independent
instrument confirmations, avoiding the circularity common in PSR-based training sets.
A Double Deep Q-Network (DQN) agent performs autonomous rover navigation, maximising
ice probability while accounting for subsurface thermal state. On a benchmark of 14
catalogued locations drawn from four published datasets, the model achieves 100\%
accuracy (14/14), with all six negative controls correctly rejected and all eight
catalogued PSRs correctly detected (F1 = 0.997, Recall = 1.000, val\_loss = 0.0109
on the validation set at 30 epochs). The platform
exposes a FastAPI backend and an interactive React web interface for point-and-click
ice probability analysis of the full 180$\times$360 lunar grid.
\end{abstract}

\vspace{0.5em}
\noindent\textbf{Keywords:} lunar water ice, permanently shadowed regions,
physics-informed neural network, deep reinforcement learning,
Lunar Reconnaissance Orbiter, subsurface thermal modelling, autonomous rover.

\hrule
\vspace{1em}

% ---------------------------------------------------------------
\section{Introduction}

The presence of water ice in lunar PSRs has been progressively confirmed by multiple
independent instruments over the past two decades. The LCROSS impact experiment
definitively detected H$_2$O in the Cabeus crater ejecta plume \citep{colaprete2010},
while Diviner temperature measurements identified cold traps capable of preserving
ice over geological timescales \citep{paige2010}. Complementary evidence was provided
by Mini-RF circular polarisation ratio (CPR) anomalies \citep{spudis2010} and LAMP
ultraviolet signatures \citep{gladstone2010}.

Despite this multi-instrument convergence, existing detection approaches typically
train models on PSR geometry alone, inadvertently reusing the same spatial prior used
to define the labels — a circularity that inflates reported performance. Furthermore,
few systems integrate subsurface thermal physics into the feature space, despite
surface temperature being a poor proxy for ice stability at depth.

This work addresses both limitations. We (1) construct labels using only
instrument-direct measurements with no overlap with the PSR spatial prior, (2) encode
subsurface temperature profiles derived from the Vasavada thermal diffusion model as
explicit input features, and (3) couple the ice detection model with a reinforcement
learning rover agent that uses subsurface state as part of its reward signal.

% ---------------------------------------------------------------
\section{Related Work}

\citet{mazarico2011} provided the foundational PSR catalogue from LOLA illumination
modelling. \citet{hayne2015} extended the Diviner analysis to a global hydration map.
\citet{williams2019} published the EPF dataset, providing direct surface-emission
temperature measurements inside PSRs. The thermal skin depth formulation used here
follows \citet{vasavada2012}, with thermal diffusivity $\kappa = 4.7 \times 10^{-7}$
m$^2$\,s$^{-1}$ appropriate for the uppermost lunar regolith.

Machine learning approaches to planetary surface analysis have been applied to Mars
and Moon imagery \citep{sato2014}, but physics-informed architectures integrating
multi-instrument LRO data with subsurface modelling have not, to our knowledge, been
previously published in open-source form.

% ---------------------------------------------------------------
\section{Data}

\subsection{Grid and Coordinate Convention}

All arrays use a global 180$\times$360 grid at 1$^\circ$/pixel. Grid indices
$(i, j)$ map to geographic coordinates as:
\begin{equation}
  \phi = i - 90, \qquad \lambda = j - 180
\end{equation}
where $\phi \in [-90^\circ, 90^\circ]$ is latitude and $\lambda \in [-180^\circ,
180^\circ]$ is longitude. This single coordinate convention is enforced across all
pipeline modules to prevent label misalignment.

\subsection{Instrument Datasets}

\begin{table}[H]
\centering
\caption{LRO instrument data integrated in the pipeline.}
\label{tab:data}
\begin{tabular}{lllr}
\toprule
Instrument & Data product & Acquisition method & Size \\
\midrule
Diviner LRO  & Polar temperature     & vsicurl GeoTIFF, $|\phi| \geq 60^\circ$ & $\sim$260 KB\\
Diviner EPF  & PSR direct temps      & PDS direct download       & $\sim$15 MB \\
LROC WAC     & 32$\times$32 patches  & vsicurl USGS COG per PSR  & $\sim$1 MB  \\
Mini-RF      & CPR polar radar       & HTTP range requests       & 59 MB       \\
LAMP         & UV H$_2$O signature   & PDS FITS download         & $\sim$10 MB \\
\bottomrule
\end{tabular}
\end{table}

Selective streaming avoids downloading full-resolution global mosaics
(typically $>$10 GB). Only polar strips ($|\phi| \geq 60^\circ$) are retrieved
from Diviner; LROC patches are fetched per-PSR via Cloud-Optimised GeoTIFF
(COG) range requests against the USGS tile server.

\subsection{Label Construction}

Labels are binary ($y \in \{0, 1\}$) and are derived from three independent
instrument sources fused with per-source confidence weights:

\begin{table}[H]
\centering
\caption{Label sources, confidence weights, and positive counts.}
\label{tab:labels}
\begin{tabular}{llrl}
\toprule
Source & Confidence & Positives & Reference \\
\midrule
LCROSS Cabeus impact    & 1.0 & 1      & \citet{colaprete2010} \\
Diviner EPF direct      & 1.0 & $\sim$6 & \citet{williams2019} \\
LAMP UV Shackleton      & 0.9 & 2      & \citet{gladstone2010} \\
Mini-RF CPR $> 1.0$     & 0.8 & 14{,}655 & \citet{spudis2010}  \\
Diviner cold trap       & 0.7 & $\sim$50 & \citet{paige2010}  \\
\midrule
\textbf{Total}          &     & \textbf{14{,}656} & \\
\bottomrule
\end{tabular}
\end{table}

PSR geometry is used only to locate known confirmed sites, never as a positive
label itself. This decouples the training signal from the spatial prior and
prevents performance inflation through circularity.

% ---------------------------------------------------------------
\section{Model Architecture}

\subsection{LunarCNN}

The detection model is a three-branch multimodal network (Listing~\ref{lst:arch}):

\paragraph{CNNEncoder.} Processes 64$\times$64 greyscale image patches through
three Conv–BN–ReLU–MaxPool blocks (channels: 1$\to$16$\to$32$\to$64),
followed by a fully-connected layer (4096$\to$128) producing a 128-dimensional feature vector.

\paragraph{PhysicsEncoder.} Receives a 5-dimensional physics feature vector:
\begin{equation}
  \mathbf{f}_p = \bigl[S_\text{norm},\ \phi_\text{norm},\ T_{0.1},\ T_{0.5},\ T_{1.0}\bigr]
\end{equation}
where $S_\text{norm} = S/1361$ is insolation normalised to the solar constant,
$\phi_\text{norm} = \phi/90$ is normalised latitude, and $T_{0.1}$, $T_{0.5}$,
$T_{1.0}$ are subsurface temperatures at 0.1, 0.5, and 1.0 m depth (normalised).
A three-layer MLP (5$\to$32$\to$64$\to$64, BatchNorm + ReLU) produces a
64-dimensional output.

\paragraph{FusionHead.} Concatenates the CNN (128-D) and physics (64-D) vectors
into a 192-D representation, then applies a linear stack
$192 \to 128 \to 64 \to 1$ with sigmoid activation to produce $\hat{p} = P(\text{ice})$.

\begin{lstlisting}[language=Python, caption={LunarCNN forward pass (simplified).}, label=lst:arch]
img_feat   = self.cnn_encoder(image)        # (B, 128)
phys_feat  = self.physics_encoder(physics)  # (B, 64)
fused      = torch.cat([img_feat, phys_feat], dim=1)  # (B, 192)
p_ice      = self.fusion_head(fused)        # (B, 1)
\end{lstlisting}

\subsection{Subsurface Thermal Model}

Subsurface temperature profiles follow the analytical solution to the one-dimensional
thermal diffusion equation for a periodically forced surface \citep{vasavada2012}:
\begin{equation}
  T(z) = \bar{T} + \left(T_s - \bar{T}\right) \exp\!\left(-\frac{z}{z_\text{skin}}\right)
\end{equation}
where $T_s$ is the instantaneous surface temperature, $\bar{T}$ is the mean annual
temperature, and the thermal skin depth is:
\begin{equation}
  z_\text{skin} = \sqrt{\frac{\kappa P}{\pi}} \approx 0.62\ \text{m}
\end{equation}
with $\kappa = 4.7 \times 10^{-7}$ m$^2$\,s$^{-1}$ and $P = 2.551 \times 10^6$ s
(lunar synodic period). Surface temperature is derived from radiative equilibrium:
\begin{equation}
  T_s = \left(\frac{(1 - A)\,S\cos\theta + \epsilon\sigma T_\infty^4}{\epsilon\sigma}\right)^{1/4}
\end{equation}
with albedo $A = 0.12$, emissivity $\epsilon = 0.95$, and $T_\infty = 3$ K.
Profiles are computed at 20 depth layers spanning 0–2 m and stored as
\texttt{temperatura\_subsolo.npy} (shape $180 \times 360 \times 20$).

\subsection{Training}

\begin{itemize}
  \item \textbf{Dataset:} 58{,}624 examples; 14{,}656 positives (25\%)
  \item \textbf{Loss:} Weighted binary cross-entropy (handles class imbalance)
  \item \textbf{Optimiser:} Adam, lr = $10^{-3}$
  \item \textbf{Schedule:} CosineAnnealingLR
  \item \textbf{Epochs:} 30 (configurable via \texttt{TRAIN\_EPOCHS})
  \item \textbf{Uncertainty:} Monte Carlo Dropout (30 forward passes) for variance estimation
\end{itemize}

% ---------------------------------------------------------------
\section{Autonomous Rover Navigation}

\subsection{Markov Decision Process Formulation}

The rover navigation task is cast as an episodic MDP with:

\paragraph{State space ($\text{obs\_dim} = 6$):}
\begin{equation}
  \mathbf{s} = \bigl[\phi_n,\ \lambda_n,\ E_n,\ S_n,\ \hat{p}_\text{ice},\ T^{(0.5\text{m})}_n\bigr]
\end{equation}
where $\phi_n$, $\lambda_n$ are normalised grid coordinates, $E_n$ is normalised
remaining energy, $S_n$ is normalised insolation, $\hat{p}_\text{ice}$ is the CNN
ice probability, and $T^{(0.5\text{m})}_n$ is the normalised 0.5 m subsurface
temperature at the current position.

\paragraph{Action space:} Four discrete moves — North, South, East, West.

\paragraph{Reward:}
\begin{equation}
  r = \Delta\hat{p}_\text{ice} - c_E\,\Delta E + r_\text{discover} + r_\text{subsurface}
\end{equation}
\begin{equation}
  r_\text{subsurface} = \max\!\bigl(0,\ 1 - T^{(0.5\text{m})}_n\bigr) \times 0.4
\end{equation}
The subsurface bonus incentivises the agent to prefer cold-at-depth regions even
when surface temperature is ambiguous — a physically motivated exploration prior.

\subsection{Double DQN Agent}

The agent uses Double Deep Q-Learning \citep{vanHasselt2016} with:
\begin{itemize}
  \item Experience replay buffer of 10{,}000 transitions
  \item Target network with soft update: $\theta^- \leftarrow \tau\theta + (1-\tau)\theta^-$, $\tau = 0.005$
  \item $\varepsilon$-greedy exploration with exponential decay over 5{,}000 steps
  \item Q-network: MLP $6 \to 128 \to 128 \to 4$
  \item Training: 500 episodes
\end{itemize}

% ---------------------------------------------------------------
\section{Results}
\label{sec:results}

\subsection{Detection Performance}

\begin{table}[H]
\centering
\caption{CNN training and validation metrics (30 epochs, CPU).}
\label{tab:results}
\begin{tabular}{lr}
\toprule
Metric & Value \\
\midrule
val\_loss & 0.0109 \\
F1 score  & 0.997  \\
Recall    & 1.000  \\
Accuracy  & 0.998  \\
\bottomrule
\end{tabular}
\end{table}

\subsection{Benchmark Against Published Catalogues}

\begin{table}[H]
\centering
\caption{Benchmark on 14 catalogued locations from four publications.
  \checkmark~= correct prediction; $\times$~= incorrect.}
\label{tab:benchmark}
\begin{tabular}{llccc}
\toprule
Location & Source & $\hat{p}$ & Expected & Result \\
\midrule
\multicolumn{5}{l}{\textit{South pole — confirmed ice}} \\
Cabeus      & \citet{colaprete2010} & 0.983 & ICE & \checkmark \\
Shackleton  & \citet{gladstone2010} & 1.000 & ICE & \checkmark \\
Haworth     & \citet{spudis2010}    & 1.000 & ICE & \checkmark \\
Nobile      & \citet{hayne2015}     & 1.000 & ICE & \checkmark \\
Amundsen    & \citet{mazarico2011}  & 1.000 & ICE & \checkmark \\
\midrule
\multicolumn{5}{l}{\textit{North pole — confirmed ice}} \\
Hermite     & \citet{paige2010}     & 1.000 & ICE & \checkmark \\
Peary       & \citet{spudis2010}    & 1.000 & ICE & \checkmark \\
Whipple     & \citet{mazarico2011}  & 1.000 & ICE & \checkmark \\
\midrule
\multicolumn{5}{l}{\textit{Negative controls — no ice}} \\
Equator 0°N & —                     & 0.000 & NO ICE & \checkmark \\
Mare Tranq. & \citet{sato2014}      & 0.000 & NO ICE & \checkmark \\
Copernicus  & —                     & 0.000 & NO ICE & \checkmark \\
Tycho       & —                     & 0.000 & NO ICE & \checkmark \\
Mid-lat 40°N & —                    & 0.000 & NO ICE & \checkmark \\
Mid-lat 40°S & —                    & 0.000 & NO ICE & \checkmark \\
\midrule
\textbf{Total} & & & & \textbf{14/14 (100\%)} \\
\bottomrule
\end{tabular}
\end{table}

\subsection{Scientific Validation}

Validation against 6 PSRs with known Diviner EPF measurements confirms that
predicted probabilities are consistent with published surface temperatures and
cold-trap stability criteria ($T < 110$ K, \citealt{paige2010}).

An earlier iteration of the label pipeline benchmarked at 12/14 (86\%), failing on
Nobile and Amundsen — both lower-confidence PSRs (conf~$\leq 0.75$) at non-extreme
south-polar longitudes ($\lambda = +53.5^\circ$ and $+84.7^\circ$) without direct
LCROSS, EPF, or LAMP confirmation, relying instead on Diviner cold-trap modelling
\citep{hayne2015} and LOLA illumination \citep{mazarico2011} alone. The cause was
traced to \texttt{generate\_labels.py} and \texttt{generate\_scientific\_data.py}
using approximate Euclidean grid-index distance rather than true haversine
great-circle distance when checking PSR proximity, which under-propagated positive
label signal near these two non-extreme-longitude sites. After correcting both
modules to use vectorised haversine distance throughout, label coverage near Nobile
and Amundsen became accurate and the current model classifies all 14 catalogued
locations correctly (Table~\ref{tab:benchmark}).

\subsection{Interpretability and Calibration}

To verify that predictions are grounded in physically meaningful signal rather
than an opaque non-linear fit, we apply three complementary interpretability
techniques (\texttt{model/interpret.py}): Grad-CAM \citep{selvaraju2017} on the
last convolutional block of the \texttt{CNNEncoder}; gradient$\times$input
attribution over the five \texttt{PhysicsEncoder} inputs, with a SHAP
GradientExplainer \citep{lundberg2017} variant that isolates the physics
branch by zeroing the image input; and a reliability diagram with Expected
Calibration Error (ECE) from Monte Carlo Dropout \citep{gal2016} (30 stochastic
forward passes).

Early development of this pipeline surfaced a bug worth reporting: the script
that drives per-coordinate Grad-CAM and attribution left \texttt{insolacao} and
\texttt{temperatura} at their function defaults (0.0 and \texttt{None}) instead
of querying the real grid values used everywhere else in the codebase, so every
probed coordinate silently received a zero-insolation, zero-subsurface-temperature
input — identical to a permanently shadowed region regardless of true location.
This inflated ice probability at all four negative controls in the interpretability
probe set ($\hat p \approx 0.96$) and collapsed attribution to $\approx$100\% on
\texttt{lat} everywhere, since the other four features were identically zero and
therefore trivially had zero gradient$\times$input. After fixing the script to
query \texttt{insolacao.npy} and derive surface temperature per coordinate
(mirroring \texttt{model/benchmark.py}), the interpretability probe set —
12 locations, a distinct but overlapping set from the 14-location benchmark of
Table~\ref{tab:benchmark} — scores 11/12, with only Sylvester (a PSR outside the
main benchmark catalogue) misclassified. Attribution is now spread across all
five features rather than collapsing onto \texttt{lat}: for positive (ice)
locations, mean attribution is lat 28.2\%, sub\_0.1m 27.2\%, sub\_0.5m 22.8\%,
sub\_1.0m 16.8\%, insolacao 5.1\%; for negative controls, sub\_1.0m 34.4\%,
sub\_0.1m 32.7\%, sub\_0.5m 23.3\%, insolacao 7.2\%, lat 2.4\%. Subsurface
thermal features dominate attribution for negatives, consistent with the model
using warm subsurface temperature — not merely latitude — to reject non-PSR
sites. Calibration on the validation set gives ECE~$=0.0905$, indicating mild
overconfidence (ECE~$>0.05$) rather than the well-calibrated regime, which we
report as an open item rather than round down to a more favourable figure.

We include this bug and its fix in the manuscript deliberately: an
interpretability pipeline is only trustworthy if its own inputs are verified,
and a plausible-looking but wrong explanation (100\% latitude attribution) is a
more dangerous failure mode than an obviously broken one, since it would have
been easy to mistake for a genuine finding about model behaviour.

% ---------------------------------------------------------------
\section{System Architecture}

The platform consists of three layers:

\paragraph{Backend (FastAPI).} Serves five endpoints with rate limiting (slowapi).
The \texttt{/analisar} endpoint returns eight fields per query:
\texttt{probabilidade\_gelo}, \texttt{variancia}, \texttt{confianca},
\texttt{temperatura}, \texttt{temperatura\_subsolo[3]}, \texttt{insolacao},
\texttt{insolacao\_atual}, \texttt{fase\_lunar}. Authentication uses
\texttt{X-API-Key} headers with CSP and HSTS headers in production.

\paragraph{Frontend (React + Leaflet).} A landing page with eight sections —
including an interactive lunar map where the user clicks to trigger ice analysis —
served via Vite at port 5173. Rover trajectories are drawn in real time as
grid-to-degree converted polylines.

\paragraph{Infrastructure.} Docker multi-stage build on \texttt{python:3.13-slim};
GitHub Actions CI; data and model volumes persist across container rebuilds.

% ---------------------------------------------------------------
\section{Discussion}

The 100\% benchmark accuracy (14/14) reflects a physics-informed label pipeline
free of circular construction rather than an aggregate figure to be taken at face
value: two of the fourteen sites (Nobile, Amundsen) are the lowest-confidence PSRs
in the catalogue (conf~$\leq 0.75$) and lack direct instrument confirmation beyond
Diviner cold-trap modelling, and only classify correctly after the haversine label
fix described in Section~\ref{sec:results} — a reminder that a perfect score on
this benchmark is sensitive to label-pipeline correctness, not solely to model
capacity. The value of the system lies not in the aggregate benchmark figure alone
but in (1) the absence of circular label construction, (2) the explicit physics
encoding of subsurface temperature, and (3) the provision of per-prediction
uncertainty via MC Dropout.

The subsurface thermal bonus in the reward function of the DQN agent constitutes
a physically motivated exploration prior. Regions cold at 0.5 m depth are
preferentially visited even when surface albedo or insolation measurements are
ambiguous — a heuristic directly grounded in ice stability theory \citep{paige2010}.

A limitation of the current label set is the dominance of Mini-RF CPR positives
(99.99\% of positives), whose confidence weight (0.8) is lower than direct
EPF measurements. Future work should prioritise integrating additional Diviner EPF
observations and LAMP count-rate maps to strengthen the high-confidence label pool.

\subsection{Out-of-Distribution Generalisation}

To probe robustness beyond the random 80/20 split used for the headline metrics,
we ran a polar-quadrant cross-validation: the model was trained on all data
excluding one polar quadrant ($|\text{lat}| > 60^\circ$, matching the polar-band
convention already used for insolation zoning) and evaluated exclusively on the
withheld quadrant, for both hemispheres independently (30 epochs each,
\texttt{model/cross\_validate.py}).

The result is a genuine negative finding. Training without any south-polar
example and testing only on that quadrant (56.6\% positive) yields
$F_1 = 0.000$ across all 30 epochs — the model detects no ice at all in a
latitude range absent from training. The mirrored north-polar fold is unstable
rather than reliably better: $F_1$ oscillates between $\approx 0.76$ and $0.000$
across epochs rather than converging monotonically, and the best reported value
($F_1 = 0.731$, recall $= 1.000$, precision $= 0.576$) corresponds to the
epoch with lowest validation loss rather than a stable operating point.

We attribute this to \texttt{lat\_norm} being one of the five \texttt{PhysicsEncoder}
input features: rather than learning purely transferable subsurface-thermal and
insolation physics, the model likely relies in part on latitude as a
near-lookup signal, which cannot extrapolate to a latitude band never observed
during training. This does \emph{not} invalidate the headline results in
Section~\ref{sec:results} — production training uses a random split that
includes examples from every latitude band, so the deployed model has not
been shown to fail in this way — but it is a real limitation of the model's
extrapolation capability that should be disclosed rather than omitted. We
report it here as an open problem for future work (e.g. ablating
\texttt{lat\_norm}, or regularising explicitly against latitude-based
shortcut learning) rather than resolving it prior to this submission.

% ---------------------------------------------------------------
\section{Conclusion}

We have described Lunar Ice Intelligence, a physics-informed deep learning platform
for lunar water ice detection. Key contributions are: a multimodal LunarCNN
integrating image and subsurface thermal features; a label pipeline free of
PSR-geometry circularity; a Double DQN rover agent with a physics-grounded
subsurface reward; and an end-to-end open-source system from raw LRO data to
interactive web interface. The software is released under the Apache License 2.0
and archived with a CITATION.cff for reproducible citation.

% ---------------------------------------------------------------
\section*{Software Availability}

Source code, trained weights (\texttt{pesos.pth}, \texttt{rl\_pesos.pth}),
and data pipeline are available under the Apache License 2.0.
Cite as: Almeida da Silva, Y. (2026). \textit{Lunar Ice Intelligence}.
ORCID: \href{https://orcid.org/0009-0007-0094-0915}{0009-0007-0094-0915}.

% ---------------------------------------------------------------
\section*{Acknowledgements}

The author acknowledges NASA's Planetary Data System (PDS) and the LRO science
teams (Diviner, Mini-RF, LROC, LAMP) for making the datasets publicly available.

% ---------------------------------------------------------------
\begin{thebibliography}{99}

\bibitem[Colaprete et al.(2010)]{colaprete2010}
Colaprete, A., et al. (2010).
Detection of Water in the LCROSS Ejecta Plume.
\textit{Science}, 330, 463--468.
\url{https://doi.org/10.1126/science.1186986}

\bibitem[Gal \& Ghahramani(2016)]{gal2016}
Gal, Y., \& Ghahramani, Z. (2016).
Dropout as a Bayesian Approximation: Representing Model Uncertainty in Deep Learning.
\textit{Proceedings of the 33rd International Conference on Machine Learning (ICML)}, 48, 1050--1059.

\bibitem[Gladstone et al.(2010)]{gladstone2010}
Gladstone, G. R., et al. (2010).
LAMP: The Lyman Alpha Mapping Project on the Lunar Reconnaissance Orbiter Mission.
\textit{Science}, 330, 472--476.
\url{https://doi.org/10.1126/science.1197874}

\bibitem[Hayne et al.(2015)]{hayne2015}
Hayne, P. O., et al. (2015).
Evidence for Exposed Water Ice in the Moon's South Polar Regions from Lunar
Reconnaissance Orbiter Ultraviolet Albedo and Temperature Measurements.
\textit{Journal of Geophysical Research: Planets}, 120, 1599--1615.
\url{https://doi.org/10.1002/2014JE004797}

\bibitem[Lundberg \& Lee(2017)]{lundberg2017}
Lundberg, S. M., \& Lee, S.-I. (2017).
A Unified Approach to Interpreting Model Predictions.
\textit{Advances in Neural Information Processing Systems (NeurIPS)}, 30, 4765--4774.

\bibitem[Mazarico et al.(2011)]{mazarico2011}
Mazarico, E., et al. (2011).
Illumination conditions of the lunar polar regions using LOLA topography.
\textit{Icarus}, 211, 1066--1081.
\url{https://doi.org/10.1016/j.icarus.2010.10.030}

\bibitem[Paige et al.(2010)]{paige2010}
Paige, D. A., et al. (2010).
Diviner Lunar Radiometer Observations of Cold Traps in the Moon's South Polar Region.
\textit{Science}, 330, 479--482.
\url{https://doi.org/10.1126/science.1187726}

\bibitem[Sato et al.(2014)]{sato2014}
Sato, H., et al. (2014).
Resolved Hapke parameter maps of the Moon.
\textit{Journal of Geophysical Research: Planets}, 119, 1775--1805.
\url{https://doi.org/10.1002/2013JE004580}

\bibitem[Selvaraju et al.(2017)]{selvaraju2017}
Selvaraju, R. R., Cogswell, M., Das, A., Vedantam, R., Parikh, D., \& Batra, D. (2017).
Grad-CAM: Visual Explanations from Deep Networks via Gradient-Based Localization.
\textit{Proceedings of the IEEE International Conference on Computer Vision (ICCV)}, 618--626.

\bibitem[Spudis et al.(2010)]{spudis2010}
Spudis, P. D., et al. (2010).
Initial results for the north pole of the Moon from Mini-RF, the Miniature Radio
Frequency Radar.
\textit{Geophysical Research Letters}, 37, L06204.
\url{https://doi.org/10.1029/2009GL042259}

\bibitem[van Hasselt et al.(2016)]{vanHasselt2016}
van Hasselt, H., Guez, A., \& Silver, D. (2016).
Deep Reinforcement Learning with Double Q-learning.
\textit{Proceedings of the AAAI Conference on Artificial Intelligence}, 30(1).

\bibitem[Vasavada et al.(2012)]{vasavada2012}
Vasavada, A. R., et al. (2012).
Lunar equatorial surface temperatures and regolith properties from the Diviner
Lunar Radiometer Experiment.
\textit{Journal of Geophysical Research: Planets}, 117, E00H18.
\url{https://doi.org/10.1029/2011JE003987}

\bibitem[Williams et al.(2019)]{williams2019}
Williams, J.-P., et al. (2019).
Diviner Lunar Radiometer Experiment eight-year special observation data record:
Reduction in data gaps and improvements in calibration, spatial sampling, and
noise.
\textit{Journal of Geophysical Research: Planets}, 124, 2078--2088.
\url{https://doi.org/10.1029/2018JE005803}

\end{thebibliography}

\end{document}
